% !PRINTED 580
\weberpage{580}
\weq{6}{(\varrho,\eta)(\varrho^2,\eta) = n, \quad \psi_1(\varrho)\psi_1(\varrho^2) = n,}
worin

\weq{7}{(\varrho,\eta) = \eta + \varrho\eta_1 + \varrho^2\eta_2, \quad (\varrho^2,\eta) = \eta + \varrho^2\eta_1 + \varrho\eta_2.}

Beispielsweise erhält man für $n = 7, 13$, wenn man die
primitiven Wurzeln 3, 2 zu Grunde legt und die Indextabellen
anwendet:

\[
n = 7
\]
\[
\begin{array}{c|c|c|c|c|c|c}
N & 1 & 2 & 3 & 4 & 5 & 6 \\
\hline
I & 0 & 2 & 1 & 4 & 5 & 3
\end{array}
\]
\[
n = 13
\]
\[
\begin{array}{c|c|c|c|c|c|c|c|c|c|c|c}
1 & 2 & 3 & 4 & 5 & 6 & 7 & 8 & 9 & 10 & 11 & 12 \\
\hline
0 & 1 & 4 & 2 & 9 & 5 & 11 & 3 & 8 & 10 & 7 & 6
\end{array}
\]

\[
\text{für } n = 7, \quad a + b\varrho = -(1 + 3\varrho),
\]
\[
\text{für } n = 13, \quad a + b\varrho = -(4 + 3\varrho).
\]

Aus diesen Formeln können wir leicht die cubische Glei\-chung herleiten, deren Wurzeln die $\eta, \eta_1, \eta_2$ sind. Sie hat wegen
(4) die Form
\weq{8}{\eta^3 + \eta^2 + \beta\eta + \gamma = 0,}

und die ganzen Zahlen $\beta, \gamma$ sind zu bestimmen.

Führen wir die Multiplication in (6) aus, so ergiebt sich
\[
\begin{aligned}
n &= \eta^2 + \eta_1^2 + \eta_2^2 - \eta\eta_1 - \eta\eta_2 - \eta_1\eta_2 \\
  &= (\eta + \eta_1 + \eta_2)^2 - 3\beta,
\end{aligned}
\]
also
\weq{9}{3\beta = -(n - 1).}

Die Ausführung der Cuben in (5) ergiebt, wenn wir für den
Augenblick
\[
s_3 = \eta^3 + \eta_1^3 + \eta_2^3, \quad s = \eta^2\eta_1 + \eta_1^2\eta_2 + \eta_2^2\eta,
\]
\[
s' = \eta_1^2\eta + \eta_2^2\eta_1 + \eta^2\eta_2
\]
setzen
\[
n\psi_1(\varrho) = s_3 - 6\gamma + 3s\varrho + 3s'\varrho^2.
\]
Aus (8) und (9) erhält man aber $s_3 = -n - 3\gamma$, also
\weq{10}{n[\psi_1(\varrho) + 1] = -9\gamma + 3s\varrho + 3s'\varrho^2,}
und ebenso
\weq{11}{n[\psi_1(\varrho^2) + 1] = -9\gamma + 3s\varrho^2 + 3s'\varrho,}
wozu noch, wenn man (4) in den Cubus erhebt,
\[
n - 1 = -9\gamma + 3s + 3s'
\]
kommt. Addirt man die drei letzten Gleichungen, so folgt
\[
n[\psi_1(\varrho) + \psi_1(\varrho^2)] + 3n - 1 = -27\gamma,
\]
%CONT>>
